\songtitle{\textchinese{夜夜尋龍} (Je je cam lung)}

\tag*{2026}\tag{cantopop}\tag{cantonese-lyrics}
% -> generos/pop.tex?
\begin{notes}
Lyrics by Claude following specific instructions on genre and narrative elements: a young dragon who changes dance floors every Friday looking for one face. Retrieved from Dvigitt's Suno page (V6-MINI, 14 November 2026).
\end{notes}
\begin{style}
Cantopop ballad, Cantonese male vocals in a low, dark lower register, warm and rounded; gated snare reverb and a bittersweet mix; mid-tempo, slightly danceable synth groove, pulsing synth bass, chiming synths and soft keys; pre-chorus lift, sparse bridge, final chorus with an upward key change.
\end{style}
\begin{lyrics}
\songpart{Verse 1}
\begin{chinese}
週五晚 鱗片染上霓虹光\\
今晚這一間 未見過你模樣\\
我收起爪牙 混入人潮之中\\
任煙霧纏繞 呼吸仍帶火光
\end{chinese}

\songpart{Pre-Chorus}
\begin{chinese}
上星期 藍調酒吧 有人話你曾經過\\
這星期 天台舞廳 剩低香水味未散
\end{chinese}

\songpart{Chorus}
\begin{chinese}
我是一條 年少的龍\\
每個週末 換一個舞池去尋蹤\\
千年火焰 燒不盡 這份空\\
誰能告訴我 你在哪個夢
\end{chinese}

\songpart{Verse 2}
\begin{chinese}
人人笑我 痴心妄想\\
龍怎可愛上 凡間的月光\\
我學人類 喝醉 學人類 跳舞\\
卻學不會 怎樣把你遺忘
\end{chinese}

\songpart{Bridge}
\begin{chinese}
鱗會褪色 爪會變鈍\\
但每個週五 心跳依然認真
\end{chinese}

\songpart{Final Chorus (key change up)}
\begin{chinese}
我是一條 年少的龍\\
就算夜夜落空 也不想放棄追風\\
若果最後 一間舞池有你在中\\
我願意 為你 化作凡人

我是一條 年少的龍\\
就算夜夜落空 也不想放棄追風\\
若果最後 一間舞池有你在中\\
收起雙翼 在你面前
\end{chinese}
\end{lyrics}

